\SourceLeaf{23}
\Article{21}

\emph{Sind $a, b, c$ usw. paarweise teilerfremd und ist ihr Produkt $abc$ usw. eine Potenz, etwa $=k^n$, so sind auch die einzelnen Zahlen $a, b, c$ usw. Potenzen derselben Art.}

Sei $a = l^\lambda m^\mu p^\pi$ usw., wobei $l, m, p$ usw. verschiedene Primzahlen bezeichnen. Nach \emph{Voraussetzung} ist keine davon ein Faktor von $b, c$ usw. Das Produkt $abc$ usw. enthält also den Faktor $l$ $\lambda$-mal, den Faktor $m$ $\mu$-mal usw. Nach dem \emph{vorigen Artikel} sind daher $\lambda, \mu, \pi$ usw. durch $n$ teilbar, und somit ist
\[
{}^{n}\!\surd a = l^{\frac{\lambda}{n}} \, m^{\frac{\mu}{n}} \, p^{\frac{\pi}{n}} \text{ usw.}
\]
eine ganze Zahl. Dasselbe gilt für die übrigen Zahlen $b, c$ usw.

Diese Tatsachen über Primzahlen mussten vorausgeschickt werden; nun wenden wir uns den Dingen zu, die unmittelbarer mit unserem Ziel zusammenhängen.

\Article{22}

\emph{Sind die Zahlen $a, b$, die beide durch eine Zahl $k$ teilbar sind, nach einem Modul $m$, der zu $k$ teilerfremd ist, kongruent, so sind auch $\frac{a}{k}$ und $\frac{b}{k}$ nach demselben Modul kongruent.}

Denn $a-b$ ist offenbar durch $k$ und ebenso durch $m$ teilbar (\emph{nach Voraussetzung}). Nach \emph{Artikel 19} ist also $\frac{a-b}{k}$ durch $m$ teilbar, das heißt $\frac{a}{k} \equiv \frac{b}{k}$ (mod. $m$).

Haben dagegen $m$ und $k$ bei sonst gleichen Voraussetzungen den größten gemeinsamen Teiler $e$, so ist $\frac{a}{k} \equiv \frac{b}{k}$ (mod. $\frac{m}{e}$). Denn $\frac{k}{e}$ und $\frac{m}{e}$ sind teilerfremd. Nun ist $a-b$ sowohl durch $k$ als auch durch $m$ teilbar, also auch $\frac{a-b}{e}$ sowohl durch $\frac{k}{e}$ als auch durch $\frac{m}{e}$ und damit durch $\frac{km}{ee}$. Das heißt, $\frac{a-b}{k}$ ist durch $\frac{m}{e}$ teilbar, oder $\frac{a}{k} \equiv \frac{b}{k}$ (mod. $\frac{m}{e}$).

\Article{23}

\emph{Ist $a$ zu $m$ teilerfremd und sind $e, f$ nach dem Modul $m$ inkongruent, so sind auch $ae, af$ nach $m$ inkongruent.}

Dies ist lediglich die Umkehrung des Satzes im vorigen Artikel.

Daraus folgt: Multipliziert man $a$ mit allen ganzen Zahlen von $0$ bis $m-1$ und führt die Produkte auf ihre kleinsten Reste nach dem Modul $m$ zurück, so sind diese Reste alle verschieden. Keiner dieser Reste ist $> m$, und ihre Anzahl beträgt $m$; ebenso viele Zahlen gibt es von $0$ bis $m-1$. Daher kann keine dieser Zahlen unter den Resten fehlen.

\SourceLeaf{24}
\Article{24}

\emph{Der Ausdruck $ax+b$}, in dem $a, b$ gegebene Zahlen und $x$ eine unbestimmte oder veränderliche Zahl bezeichnen, \emph{kann nach einem Modul $m$, der zu $a$ teilerfremd ist, jeder gegebenen Zahl kongruent gemacht werden.}

Sei $c$ die Zahl, zu der der Ausdruck kongruent sein soll, und sei der kleinste positive Rest von $c-b$ nach dem Modul $m$ gleich $e$. Nach dem vorigen Artikel gibt es notwendig einen Wert von $x < m$, für den der kleinste Rest des Produkts $ax$ nach dem Modul $m$ gleich $e$ wird. Sei $v$ dieser Wert. Dann ist $av \equiv e \equiv c-b$, also $av+b \equiv c$ (mod. $m$). \emph{Q. E. F.}

\Article{25}

Einen Ausdruck, der zwei Größen als kongruent angibt, nennen wir in Anlehnung an Gleichungen eine \emph{Kongruenz}. Enthält er eine Unbekannte, so heißt die Kongruenz \emph{gelöst}, wenn ein Wert der Unbekannten gefunden ist, der die Kongruenz erfüllt (eine \emph{Wurzel}). Damit ist auch klar, was eine \emph{lösbare Kongruenz und eine unlösbare Kongruenz} ist. Schließlich sieht man leicht, dass hier ähnliche Unterscheidungen wie bei Gleichungen möglich sind. Beispiele \emph{transzendenter} Kongruenzen werden später vorkommen. \emph{Algebraische} Kongruenzen werden nach dem höchsten Exponenten der Unbekannten in solche ersten, zweiten und höheren \emph{Grades} eingeteilt. Ebenso können mehrere Kongruenzen mit mehreren Unbekannten vorliegen, deren \emph{Elimination} zu untersuchen ist.

\Needspace{6\baselineskip}
\begin{center}
\emph{Lösung der Kongruenzen ersten Grades.}
\end{center}

\Article{26}

Nach Artikel 24 ist also die Kongruenz ersten Grades $ax+b \equiv c$ stets lösbar, wenn der Modul zu $a$ teilerfremd ist. Ist $v$ ein geeigneter Wert von $x$, also eine Wurzel der Kongruenz, so sind offenbar auch alle Zahlen Wurzeln, die zu $v$ nach dem Modul der gegebenen Kongruenz kongruent sind (Artikel 9). Ebenso leicht sieht man, dass alle Wurzeln zu $v$ kongruent sein müssen: Ist $t$ eine weitere Wurzel, so ist $av+b \equiv at+b$, also $av \equiv at$ und damit $v \equiv t$ (Artikel 22). Folglich gibt die Kongruenz $x \equiv v$ (mod. $m$) die vollständige Lösung der Kongruenz $ax+b \equiv c$ an.

Da Lösungen einer Kongruenz durch kongruente Werte von $x$ unmittelbar auf der Hand liegen und kongruente Zahlen in dieser Hinsicht als gleichwertig anzusehen sind, zählen wir solche Lösungen als ein und dieselbe. Da nun
\SourceLeaf{25}
unsere Kongruenz $ax+b \equiv c$ keine anderen Lösungen zulässt, sagen wir, sie sei nur auf eine Weise lösbar oder habe nur eine Wurzel. \emph{Zum Beispiel} lässt die Kongruenz $6x+5 \equiv 13$ (mod. $11$) nur solche Wurzeln zu, die $\equiv 5$ (mod. $11$) sind. Bei Kongruenzen höheren Grades ist dies anders, ebenso bei Kongruenzen ersten Grades, wenn die Unbekannte mit einer Zahl multipliziert wird, zu der der Modul nicht teilerfremd ist.

\Needspace{25\baselineskip}
\Article{27}

Es bleibt noch etwas darüber zu sagen, wie man die Lösung einer solchen Kongruenz findet. Zunächst hängt eine Kongruenz der Form $ax+t \equiv u$, deren Modul wir als teilerfremd zu $a$ voraussetzen, von der Kongruenz $ax \equiv \pm 1$ ab: Erfüllt $x \equiv r$ die letztere, so erfüllt $x \equiv \pm(u-t)r$ die erstere. Die Kongruenz $ax \equiv \pm 1$, deren Modul wir mit $b$ bezeichnen, ist gleichwertig mit der unbestimmten Gleichung $ax = by \pm 1$. Wie man diese löst, ist inzwischen gut bekannt; daher genügt es, hier den Rechenalgorithmus anzugeben.

Hängen die Größen $A, B, C, D, E$ usw. von $\alpha, \beta, \gamma, \delta$ usw. durch die Beziehungen
\[
A = \alpha, \quad B = \beta A + 1, \quad C = \gamma B + A, \quad D = \delta C + B, \quad
E = \varepsilon D + C \text{ usw.}
\]
ab, so bezeichnen wir sie der Kürze halber folgendermaßen:
\begingroup
\renewcommand{\thefootnote}{*)}
\[
A = [\alpha], \quad B = [\alpha, \beta], \quad C = [\alpha, \beta, \gamma], \quad
D = [\alpha, \beta, \gamma, \delta] \text{ usw.}\footnotemark.
\]
\footnotetext{
Diese Beziehung lässt sich viel allgemeiner betrachten; vielleicht werden wir dies bei einer anderen Gelegenheit aufgreifen. Hier fügen wir nur zwei Sätze hinzu, die bei der vorliegenden Untersuchung nützlich sind:

\[1^{\mathrm{o}}.\quad [\alpha,\beta,\gamma\mathinner{\ldotp\ldotp\ldotp\ldotp} \lambda,\mu].[\beta,\gamma\mathinner{\ldotp\ldotp\ldotp\ldotp} \lambda] -
[\alpha,\beta,\gamma\ldots \lambda][\beta,\gamma\ldots \lambda,\mu] = \pm 1,\]
Dabei ist das obere Zeichen zu nehmen, wenn die Anzahl der Zahlen $\alpha,\beta,\gamma\ldots \lambda,\mu$ gerade ist, das untere, wenn sie ungerade ist.

$2^{\mathrm{o}}$. Die Reihenfolge der Zahlen $\alpha,\beta,\gamma$ usw. kann umgekehrt werden:
\[
[\alpha,\beta,\gamma\ldots \lambda,\mu]=[\mu,\lambda\mathinner{\ldotp\ldotp\ldotp\ldotp}\gamma,\beta,\alpha].
\]

Die Beweise, die nicht schwierig sind, lassen wir hier weg.
}
\endgroup
Gegeben sei nun die unbestimmte Gleichung $ax = by \pm 1$, in der $a, b$ positiv sind. Wir dürfen annehmen, dass $a$ nicht $< b$ ist. Nach dem bekannten Algorithmus zur Bestimmung des größten gemeinsamen Teilers zweier Zahlen bilden wir durch gewöhnliche Division die Gleichungen
\[
a = \alpha b + c, \quad b = \beta c + d, \quad c = \gamma d + e \text{ usw.}
\]
Dabei sind $\alpha, \beta, \gamma$ usw. sowie $c, d, e$ usw. positive ganze Zahlen, und $b, c, d, e$ nehmen fortlaufend ab, bis wir zu $m = \mu n + 1$ gelangen,
\SourceLeaf{26}
was bekanntlich schließlich eintreten muss. Es ist also
\[
a=[n,\mu,\mathinner{\ldotp\ldotp\ldotp\ldotp} \gamma,\beta,\alpha], \quad b=[n,\mu\mathinner{\ldotp\ldotp\ldotp\ldotp} \gamma,\beta]
\]
Setzen wir nun
\[
x=[\mu,\mathinner{\ldotp\ldotp\ldotp\ldotp} \gamma,\beta], \quad y=[\mu,\mathinner{\ldotp\ldotp\ldotp\ldotp} \gamma,\beta,\alpha]
\]
so ist $ax=by+1$, wenn die Anzahl der Zahlen $\alpha,\beta,\gamma\mathinner{\ldotp\ldotp\ldotp\ldotp} \mu,n$ gerade ist, und $ax=by-1$, wenn sie ungerade ist. \emph{Q. E. F.}

\Article{28}

Der berühmte Euler gab zuerst die allgemeine Lösung solcher unbestimmten Gleichungen an, in \emph{Comment. Petrop. T.} VII, \emph{S.} 46. Seine Methode besteht darin, andere Unbekannte an die Stelle von $x, y$ zu setzen, und ist inzwischen gut bekannt. Der berühmte La Grange ging etwas anders vor. Aus der Theorie der Kettenbrüche folgt nämlich: Verwandelt man den Bruch $\frac{b}{a}$ in den Kettenbruch
\[
\cfrac{1}{\alpha + \cfrac{1}{\beta + \cfrac{1}{\gamma + \text{usw.} + \cfrac{1}{\mu + \frac{x}{n}}}}}
\]
und wandelt diesen nach Weglassen seines letzten Teils $\frac{x}{n}$ wieder in einen gewöhnlichen Bruch $\frac{x}{y}$ um, so ist $ax=by\pm1$, sofern $a$ zu $b$ teilerfremd ist. Beide Methoden führen zum selben Algorithmus. La Granges Untersuchungen finden sich in \emph{Hist. de l'Ac. de Berlin Année} 1767, \emph{S.} 175, und zusammen mit weiteren Untersuchungen in den \emph{Zusätzen zur französischen Übersetzung von Eulers Algebra}.

\Article{29}

Die Kongruenz $ax+t \equiv u$, deren Modul nicht zu $a$ teilerfremd ist, lässt sich leicht auf den vorigen Fall zurückführen. Sei $m$ der Modul und der größte gemeinsame Teiler von $a, m$ gleich $\delta$. Zunächst erfüllt offenbar jeder Wert von $x$, der die Kongruenz nach dem Modul $m$ erfüllt, sie auch nach dem Modul $\delta$ (Artikel 5). Es ist aber stets $ax \equiv 0$ (mod. $\delta$), da $\delta$ ein Teiler von $a$ ist. Die gegebene Kongruenz ist daher unlösbar, wenn nicht $t \equiv u$ (mod. $\delta$), das heißt, wenn $t-u$ nicht durch $\delta$ teilbar ist.

\SourceLeaf{27}
Setzen wir also $a = \delta e$, $m = \delta f$, $t - u = \delta k$, so ist $e$ zu $f$ teilerfremd. Dann ist die gegebene Kongruenz $\delta e x + \delta k \equiv 0 \ (\mathrm{mod.}\ \delta f)$ gleichwertig mit $ex + k \equiv 0 \ (\mathrm{mod.}\ f)$. Das heißt, jeder Wert von $x$, der die letztere erfüllt, erfüllt auch die erstere, und umgekehrt. Denn $ex+k$ ist offenbar durch $f$ teilbar, wenn $\delta e x + \delta k$ durch $\delta f$ teilbar ist, und umgekehrt. Wie man die Kongruenz $ex+k\equiv 0\ (\mathrm{mod.}\ f)$ löst, haben wir oben gezeigt. Daraus folgt zugleich: Ist $v$ einer der Werte von $x$, so gibt $x \equiv v \ (\mathrm{mod.}\ f)$ die vollständige Lösung der gegebenen Kongruenz an.
